\section{Conclusion}
\label{sec:conclusion}

We introduced a cheap, interventional, success-conditioned reach readout that asks a
question success rate cannot: when a VLA policy appears to generalize spatially, including
on the rollouts it gets right, is the reach grounded to the displaced object or anchored to
the canonical training location and tolerated by a wide grasp basin? Using it on a
competent SmolVLA across two LIBERO suites, we found that the same counterfactual
fine-tuning basin-widens on a wide-basin suite and reaches object-near on a tight-basin
suite---opposite mechanisms behind identical-looking success gains. A boundary analysis
showed the readout also rejects tempting mechanistic explanations that fail to hold across
suites, underscoring that the robust contribution is the instrument, not any single
mechanism.

\paragraph{Limitations.}
Our study uses a single policy family (SmolVLA) and simulation (LIBERO). The readout is
model-agnostic by construction---it needs only object displacement at reset and endpoint
logging---but a second-policy datapoint is empirically pending: our attempts to obtain one
(ACT and a pi0.5 checkpoint) were blocked by policy-integration and compute-contention
issues, not by anything intrinsic to the readout.
%% <!-- DATA_NEEDED: GEN_2ND_POLICY — one second-policy base-only readout (u_succ / rho_succ / displaced success, one suite, 50 mm) to substantiate model-agnostic reuse -->
Seed budgets are set by a single-GPU regime; the mechanistic boundary results in
\Cref{sec:boundary} are correspondingly underpowered at the task-cluster level and are
reported as such.

\paragraph{Future work.}
The readout is released as a drop-in tool; the immediate next step is to apply it to
additional policy families and to real-robot rollouts, where basin geometry and grounding
may diverge further. Contact- and lift-time endpoints suggest a finer question---when and
where an anchored plan is corrected---that the readout is already positioned to measure.
